% Problem 2-1
\begin{itemize}
  \item[\textbf{a.}]
    \begin{proof}
      Since \textsc{Insertion-Sort} has a running time of $\Theta(n^2)$,
      due to the nested \textbf{while} loop of lines 5-7 (p. 26), and the
      loop of lines 1-7, if we are instead sorting $n$ items divided into
      $n/k$ sublists, the outer \textbf{for} loop will run $k$ times for
      each sublist. Since each sublist will contain $k$ items, the
      \textbf{while}
      loop will run $k - 1$ times for each sublist, so we have a running time
      of $\Theta(k^2)$ for each sublist, and since there are $n/k$ sublists,
      the total running time will be
      \[
        \Theta\biggl(\frac{n}{k}\cdot k^2\biggr) = \Theta(nk).
      \]
    \end{proof}
  \item[\textbf{b.}] If we merge the sublists when their length
    reaches $k$, we are
    removing $\lg{k}$ levels from the recursion tree on p. 38, so we have
    \[
      \lg{n} - \lg{k} = \lg{n/k}
    \]
    levels on the tree. Now, from \textbf{a.}, we have the
    \textsc{Insertion-Sort}
    algorithm on the $n/k$ sublists with a running time of $\Theta(nk)$, so the
    new algorithm will have a running time of
    \[
      \Theta(nk + n\lg{(n/k)}).
    \]
    Note that the two terms cannot be combined: neither dominates the other for
    all $k$, since $nk$ is the larger when $k$ is close to $n$ and
    $n\lg{(n/k)}$ is the larger when $k$ is constant.
  \item[\textbf{c.}] We want the largest $k$ for which
    \[
      \Theta(nk + n\lg{(n/k)}) = \Theta(n\lg{n}),
    \]
    and expanding $\lg{(n/k)} = \lg{n} - \lg{k}$ gives
    \[
      nk + n\lg{n} - n\lg{k}.
    \]
    The term $n\lg{n} - n\lg{k}$ never exceeds $n\lg{n}$, so the sum is
    $O(n\lg{n})$ precisely when $nk = O(n\lg{n})$, i.e. when $k = O(\lg{n})$.
    Conversely, if $k = O(\lg{n})$ then $\lg{k} = O(\lg\lg{n})$ and
    \[
      nk + n\lg{n} - n\lg{k} \geq n\lg{n} - n\lg\lg{n} = \Omega(n\lg{n}),
    \]
    so the running time is $\Theta(n\lg{n})$ exactly when $k = \Theta(\lg{n})$,
    and this is the largest such $k$ asymptotically. Any $k = \omega(\lg{n})$
    makes the $nk$ term $\omega(n\lg{n})$, so the modified algorithm would then
    be asymptotically slower than standard merge sort.
  \item[\textbf{d.}] Asymptotically, part \textbf{c.} permits any
    $k = \Theta(\lg{n})$. In practice the constants hidden by the
    $\Theta$-notation are what matter, so $k$ should be chosen as the largest
    sublist length at which \textsc{Insertion-Sort} is still measurably faster
    than \textsc{Merge-Sort} on the target machine, which is determined
    empirically rather than analytically.
\end{itemize}
